\documentclass{article}

\usepackage[utf8]{inputenc}
\usepackage[T1]{fontenc}
\usepackage[english]{babel}
\usepackage{microtype}

\usepackage{amsmath, amssymb, amsthm}
\usepackage{mathtools}
\usepackage{totcount}
\usepackage{enumitem}

\usepackage{newtxtext,newtxmath}

\usepackage{titlesec}
\titleformat{\chapter}[display]
  {\normalfont\huge\bfseries} % Format of the text
  {Lecture \thechapter}       % Label (e.g., "Lecture 1")
  {0.5em}                     % Spacing between label and title
  {\Huge}                     % Format of the title

% Theorem-like environments
\newtheorem{theorem}{Theorem}[section]
\newtheorem{lemma}[theorem]{Lemma}
\newtheorem{proposition}[theorem]{Proposition}
\newtheorem{corollary}[theorem]{Corollary}

\theoremstyle{definition}
\newtheorem{definition}[theorem]{Definition}
\newtheorem{example}[theorem]{Example}

\theoremstyle{remark}
\newtheorem{remark}[theorem]{Remark}

\title{Practice-Final Exam for Elementary Number Theory}
\author{Egbert Rijke}
\date{December 4th, 2025}

\usepackage{hyperref}
\hypersetup{
    colorlinks=true,
    linkcolor=purple,
    filecolor=magenta,
    urlcolor=blue,
    citecolor=purple,
    pdftitle={Lecture notes on Number Theory},
    pdfpagemode=FullScreen,
}
\usepackage[capitalize,english]{cleveref}

  % We set up the exercise environment, which produces list environment in a new
% unnumbered subsection that also gets mentioned in the table of contents.

\newtotcounter{exercisecounter}
\newcommand{\exercise}{\item\stepcounter{exercisecounter}}

\newlist{exenum}{enumerate}{1}
\setlist[exenum]{noitemsep,label=\arabic*}
  %,ref=\thechapter.\arabic*}

\crefname{exenumi}{Exercise}{Exercises}
  
\newlist{subexenum}{enumerate}{1}
\setlist[subexenum]{noitemsep,label=(\alph*),ref=\theexenumi.\alph*}
\crefname{subexenumi}{Exercise}{Exercises}
  
\newenvironment{exercises}
{%
\subsection*{Exercises}%
\addcontentsline{toc}{section}{Exercises}%
\sectionmark{Exercises}%
\begin{exenum}}
{%
%\addtocounter{exercisecounter}{\theexenumi}%
\end{exenum}}

\newcommand{\legendre}[2]{%
  \mathchoice
    {\left(\frac{#1}{#2}\right)}% display style
    {\bigl(\tfrac{#1}{#2}\bigr)}% text style
    {\bigl(\tfrac{#1}{#2}\bigr)}% script style
    {\bigl(\tfrac{#1}{#2}\bigr)}% scriptscript style
}

\begin{document}
\maketitle

\begin{exercises}
  \exercise Solve $x^{86}\equiv 7\pmod{29}$. \\
  Hint: Use Theorem 13.1.1.
  \exercise Find two different positive integers $n$ for which $\phi(n)=20$. \\
  Hint: Use Theorem 12.2.5.
  \exercise Determine whether the congruence $x^2+14x-35\equiv 0\pmod{57}$ has a solution. \\
  Hint: Use Theorem 18.1.2.
  \exercise Compute the Legendre symbols
  \begin{equation*}
    \legendre{337}{373}\qquad\legendre{373}{733}\qquad\text{and}\qquad\legendre{733}{337}.
  \end{equation*}
  Hint: Use Theorem 19.2.1.
  \exercise Determine for which primes $p\pmod{20}$ we have $\legendre{-5}{p}=1$. \\
  Hint: Use Theorem 19.2.3.
  \exercise
  \begin{subexenum}
  \item Find the least primitive root $g$ modulo $19$, and give a table of the congruence classes of its powers.
  \item Find all primitive roots modulo $19$.
  \item For each $k\mid\phi(19)$, find all elements of order $k$ modulo $19$.
  \item For each integer $1,\ldots,18$, find its index base $g$ modulo $19$.
  \item Solve the linear congruence
    \begin{equation*}
      7x\equiv 8\pmod{19}
    \end{equation*}
    using indices.
  \end{subexenum}
\end{exercises}

\end{document}
